% ===== 原文 pp.36--38 =====
% =====================================================================
%  tr06a.tex —— §6「未来方向」前半（原文 pp.36--38）
%  负责人：tr-06
%  覆盖：原文 p.36 的 §6 标题起，到 p.38 末段那句「…called the $T_G$ theory, or」
%        为止（该句跨页到 p.39 第一行；按 STYLE_GUIDE ⑧「开始者写完」，本片补完
%        整个句子，并顺势写到 p.39 第一段 P13 的段落边界为止，以免在段落中间
%        切分文件、造出原文没有的分段）。故 p.39 的页码注释出现在本片末尾；
%        p.39 第二段起的正文见 tr06b.tex。
%  本块无编号公式（§6 全文没有 \tag）、无图、无表。
% =====================================================================

% [原文 p.36]
\section{未来方向}

至此，我们已经走过了 2d/4d 对应关系最基础的部分：我们从 $\SU(2)$ 型的六维
$\cN{=}(2,0)$ 理论出发。它在带穿刺点的 Riemann 曲面上的紧致化，可以用 $\SU(2)$
矢量多重态与三基表示半超多重态给出拉格朗日描述。于是我们可以计算它在
$S^1\times S^3$ 上或在 $S^4$ 上的配分函数。我们看到，该配分函数分别等于二维
$q$-形变 $\SU(2)$ 杨--米尔斯理论的配分函数，或 Liouville 理论的配分函数。现在，
我们有了一个途径：通过直接研究 $S^1\times S^3$ 上或 $S^4$ 上的六维理论来理解
这一结果。

在这个最基本的情形里，我们的理解仍存在空白；而且还有许多推广的途径。我们就以
梳理这些要点来结束本综述。若有读者参与进来、解决这里提到的某些问题，那将令
作者十分高兴；若有人能开辟出这里甚至未曾提及的新方向，那就更加妙不可言。

% [Tutor Note: p.36，§6 开头这两段是全书内容的「一句话路线图」：六维理论 →
%  在 Riemann 曲面 $C_2$ 上紧致化得到 4d class S 理论 → 在 $X_4$ 上算配分函数
%  → 等于 2d 的 $q$-形变杨--米尔斯理论 / Liouville 理论。四步之中，六维理论本身
%  （第一步与第四步）没有拉格朗日量，而 4d 理论（第二、三步）有。本节所有的
%  「不知道」几乎都源于这个「半边有拉格朗日量、半边没有」的结构；作者在 p.39
%  把它命名为「困难的守恒定律」（conservation law of the difficulty）。
%  记住这条主线，四个小节就很好定位：6.1 是 SU(2) 情形下仍缺的两块拼图；
%  6.2 是同一套构造里还能放入哪些「对象」（算符）；6.3 是把型 SU(2) 换成任意
%  ADE 型；6.4 是把 $X_4$ 换成别的时空、或把六维理论换成别的理论。]

\subsection{SU(2) 情形下的两个未解决问题}

在这个最基本的情形里，本综述中那些只是被断言成立的数学关系，大多甚至已经被
严格证明。在 Liouville 理论这一侧，我们理解上的唯一一大空白在于一般 $\SU(2)$
型 class S 理论在 $S^4$ 上的局域化：我们仍然不知道如何从处于 $\SU(2)^3$ 三基
表示中的一个真正的%
% [原文 p.37]
半超多重态得到瞬子贡献，参见例如 [38]。由于这一问题，$S^4$ 配分函数只能对
线性或环形的箭图计算。作者希望有感兴趣的读者能在这点上找到推进的
办法。

% [Tutor Note: p.36--37，第一个具体未解决问题「how to obtain the instanton
%  contribution from a genuine half-hypermultiplet」卡在哪里：4d 的瞬子贡献
%  一向靠 ADHM 构造 / Nekrasov 的局域化公式来算，而那套办法要求物质场是基础
%  表示或双基表示（每个物质场只耦合两个规范群）。三穿刺球面对应的物质是
%  「$\SU(2)^3$ 三基表示的半超多重态」，它同时耦合三个 $\SU(2)$，不是任何
%  标准 quiver 的物质场，因此没有已知的 ADHM 构造。这正是作者接下来说
%  「$S^4$ 配分函数只能对线性或环形 quiver 计算」的原因：只有当裤子分解
%  （pants decomposition）图是一条链或一个圈（每个规范节点最多接两个三穿刺
%  球面）时，三基表示半超多重态才能重新组合成通常的双基表示，瞬子计数才有
%  定义。反过来，§5.2 的等式（瞬子配分函数 = Liouville 共形块）提供了另一条
%  路：绕开「直接做瞬子计数」这一步，改从共形块一侧去认；一般的 3 价或更高价
%  图（即「星形」分解）在本文写作时仍是缺口。]

在 $S^1\times S^3$ 配分函数这一侧，我们只讨论了含三个参数 $(p,q,t)$ 的一般
超共形指标中的一个单参数切片。在双参数切片 $(p=0,q,t)$ 中，已知
$S^1\times S^3$ 配分函数给出 $(q,t)$-形变杨--米尔斯理论：其中杨--米尔斯
配分函数里的特征标 $\chi_\lambda(a)$ 被 Macdonald 多项式 $P_\lambda(a|q,t)$
替换。在最一般的三参数情形，四点函数的基本结合性已经在数学上得到证明，但并不
是用我们这种把无穷乘积改写为无穷和的办法，而是在 Macdonald 多项式的进一步推广
$P_\lambda(a|p,q,t)$ 之上完成的，参见例如 [39]。

% [Tutor Note: p.37，$(p,q,t)$ 三个参数的地位：4d $\cN{=}2$ 超共形指标有三个
%  fugacity 参数。本文 §5.1 实际只用了「单参数切片」——本页的叙述反推出，那相
%  当于取 $p=0$ 且 $t=q$：此时 Macdonald 多项式 $P_\lambda(a|q,t)$ 退化回普通
%  特征标 $\chi_\lambda(a)$，剩下的对应关系就是 §5.1 的 $q$-形变杨--米尔斯。
%  本节说的「双参数切片」是 $p=0$ 但保留 $t\neq q$；「三参数情形」连 $p$ 也不
%  取零。容易误会的是：$p,q,t$ 不是三个独立耦合常数，而是指标
%  （$S^1\times S^3$ 配分函数）的三个 fugacity；所谓「切片」只是看指标的某个
%  特化极限。]
% [Tutor Note: p.37，「but not as we did by rewriting the infinite product into
%  an infinite sum」指的是 §5.1 的核心技巧：(5.1) 的半超多重态贡献本是一个
%  无穷乘积，利用 (5.2) 把它改写为对表示 $\lambda$ 的无穷和，从而认出它正是
%  2d 三穿刺球面振幅。作者在这里说：一般三参数的结合性虽然是数学定理，但证明
%  不是走这条「乘积 ↔ 和」的路，而是直接在 $P_\lambda(a|p,q,t)$ 上做的。]

\subsection{SU(2) 情形下的更多对象}

到目前为止，在本综述中我们只用到了一种特定的余维 2 算符，即 $\SU(2)$ 型六维
$\cN{=}(2,0)$ 理论的余维 2 算符。它保持那个把算符世界体积映到自身的超共形
对称性；这样的算符通常称为规则的（regular）。在带规则穿刺点的 Riemann 曲面
上的紧致化给出四维超共形理论，而此前我们讨论的一直只是这一类。

也存在不规则的（irregular）余维 2 算符，它们会破坏超共形对称性。带不规则
穿刺点的紧致化可以给出四维的渐近自由理论，也可以给出一种相当不同类型的超共形
理论，即所谓 Argyres--Douglas 理论。渐近自由理论的 $S^4$ 配分函数可以写下来，
但 Argyres--Douglas 理论的 $S^4$ 配分函数从 4d 的角度看还没有被很好地理解。
与之对应的 Liouville 关联函数正在被研究，参见例如 [40,41]。2015 年，对
Argyres--Douglas 理论的 $S^1\times S^3$ 配分函数的理解取得了进展，参见例如
[42,43]。

% [Tutor Note: p.37，「规则/不规则穿刺点」的区别值得写清：规则穿刺点的定义
%  要求该余维 2 算符在穿刺点附近保持超共形对称性——在 Hitchin 系统的描述里，
%  规则穿刺点对应 Higgs 场在穿刺点处只有简单极点（一阶极点），不规则穿刺点
%  则允许二阶及更高阶的极点；后者会破坏超共形对称性。4d 一侧
%  要么只得到渐近自由理论（有拉格朗日量，$S^4$ 配分函数可写出），要么得到没有
%  边际耦合的 Argyres--Douglas 型 SCFT。卡点在后者：没有拉格朗日量，局域化不能
%  直接使用，所以作者说其 $S^4$ 配分函数「从 4d 的观点看」还没被理解；而在 2d
%  一侧它对应的是「不规则共形块」（irregular conformal block），不是普通共形块
%  里的初级场插入——这就是「corresponding Liouville correlators are being
%  explored」这句话的含义。$S^1\times S^3$ 情形（指标）之所以 2015 年先有进展，
%  是因为指标不需要 ADHM 式的瞬子构造。]

$\SU(2)$ 型的六维 $\cN{=}(2,0)$ 理论还有余维 4 算符，它们由 $\SU(2)$ 的不可约
表示标记。由平庸的一维表示所标记的那些是平庸的，而基本的非平庸者由二重态表示
标记。回想 2d/4d 对应的基本设定：我们有定义在 $X_4\times C_2$ 上的六维理论。
由于余维 4 算符具有二维世界体积，粗略地说，我们可以有以下三种情形：

\begin{itemize}
\item 该算符在 $X_4$ 中沿两个方向延伸，并位于 $C_2$ 中的一点。这给出四维理论
  中所谓的面算符。其规范理论描述如今已被很好地理解。在 Liouville 理论中，
  这是一次退化算符（degenerate operator）的插入。它在 $q$-形变杨--米尔斯
  中的体现已在 [44] 中研究过。
\item 该算符在 $X_4$ 中沿一个方向延伸，并在 $C_2$ 中沿一个方向延伸。这给出
  四维理论中的线算符。事实上，$\SU(2)$ 型 class S 理论的所有可能的线算符都
  可以用这种方式漂亮地描述。它们也给出 Liouville 理论中的线算符 [45,46]。
% [原文 p.38]
\item 该算符位于 $X_4$ 中的一点，并覆盖整个 $C_2$。这只不过给出四维理论的
  一个点算符。当 $C_2$ 是三穿刺球面时，这个算符就是三基表示算符本身。
\end{itemize}

这些看起来或多或少已被理解；当我们从 $\SU(2)$ 推广到更大的东西时，它们会更
有意思。

% [Tutor Note: p.37--38，余维 4 算符的三个情形其实就是「2 维世界体积如何分给
%  $X_4$ 与 $C_2$」的三种分法：$2+0$、$1+1$、$0+2$，分别得到 4d 的面算符、
%  线算符、点算符（世界体积维数 $2+0=2$、$1+1=2$、$0+2=2$ 都等于余维 4 的
%  补维数）。最后一条尤其值得对照 §3：三穿刺球面 = $T_2$ 理论，其物质正是
%  $\SU(2)^3$ 三基表示 (半)超多重态；所以「卷绕整个 $C_2$ 的余维 4 算符」
%  给出这个三基表示算符本身，这句话说的是：三穿刺球面的物质就是六维理论的一个
%  余维 4 算符。]
% [Tutor Note: p.38，中间那条「线算符」：对 $\SU(2)$ 的 class S 理论，线算符
%  按电磁荷 $(m,n)\in\Z^2$ 分类（Dirac 量子化 + 对偶群的作用），这里作者说它们
%  都能由「在 $C_2$ 中缠绕一个 1-闭链的余维 4 算符」给出：双重态给出算符，
%  不同的 1-闭链给出不同的线算符。这一几何图像的好处是把 4d 的 Wilson--'t
%  Hooft 线算符翻译成了 6d 的算子谱问题。]

余维 2 算符也可以用多种方式放进这个框架：

\begin{itemize}
\item 该算符在 $X_4$ 中沿四个方向延伸，并位于 $C_2$ 中的一点。这就是我们一直
  在谈论的穿刺点，它会改变四维理论。
\item 该算符在 $X_4$ 中沿三个方向延伸，并在 $C_2$ 中缠绕一条线。这本应在
  class S 理论中给出一个畴壁算符，但似乎还没有被研究过。
\item 该算符在 $X_4$ 中沿两个方向延伸，并缠绕整个 $C_2$。这给出四维理论中的
  一个面算符，并改变 $C_2$ 上的理论。当 $X_4=S^4$ 时，$C_2$ 上的理论是
  $\mathrm{SL}(2)$ WZW 理论 [47]；其瞬子贡献此前已被研究过 [48]。
\end{itemize}

% [Tutor Note: p.38，第二个 bullet 是本节第二个「空白」：缠绕 $C_2$ 中一条线、
%  在 $X_4$ 中延伸 3 维的余维 2 算符，本应给出 4d class S 理论里的畴壁算符
%  （domain wall，在 4d 中就是 3 维界面，把两种不同的真空或两种不同的 4d 理论
%  隔开）。作者
%  明说它「does not appear to have been studied」——这是全文列出的、最干净
%  的一个尚未开工的问题：连对应的 2d 对象都不清楚。
%  第三个 bullet 可与 §5.2 对照：$X_4=S^4$ 且无面算符时，$C_2$ 上的 2d
%  理论是 Liouville；插入这种面算符后 $C_2$ 上的理论变成 $\mathrm{SL}(2)$ WZW。二者
%  并不矛盾——Liouville 理论正是 $\mathrm{SL}(2)$ WZW 经量子 Drinfeld--Sokolov
%  （哈密顿）约化得到的，所以这里是同一 2d 物理在「约化前/约化后」的两种
%  描述。读者若只记住「$C_2$ 上住着一个 2d CFT」，这条 bullet 就容易读错。]

\subsection{利用其它类型的六维 $\cN{=}(2,0)$ 理论}

到目前为止，我们只讨论了从 $\SU(2)$ 型的六维 $\cN{=}(2,0)$ 理论出发的情形。
六维 $\cN{=}(2,0)$ 理论被认为落入 ADE 分类，即存在 $G=A_n$、$D_n$ 与
$E_{n=6,7,8}$ 型的理论。我们可以把这些理论放到 $X_4\times C_2$ 上，并在求出
所得到的 2d/4d 对应关系的过程中获得许多乐趣。

在框架中没有任何额外对象时，$\cS_G(S^1\times S^3)$ 给出规范群为 $G$
的二维 $q$-形变杨--米尔斯理论，或它的 $(q,t)$ 与 $(p,q,t)$ 推广。
$\cS_G(S^4)$ 给出 Liouville 理论的一个自然推广，称为 $G$ 型的二维
Toda 理论；它具有一个称为 $W_G$ 代数的二维对称性，而 $W_G$ 代数是 Virasoro
代数的推广。

% [Tutor Note: p.38，记号 $\cS_G(X_4)$ 是前文 $\cS(X_4)$ 的
%  推广：它指「$G$ 型六维 $\cN{=}(2,0)$ 理论在 $X_4\times C_2$ 上紧致化所得的
%  4d 理论」，下标 $G$ 标明所用的六维理论类型。§5.3.2 把结论写成
%  $\cS(X_4)=\cdots$ 的形式，看起来像是「配分函数的等式」；严格说
%  $\cS(X_4)$ 指 4d 理论本身，它的配分函数才等于右边那个 2d 理论的
%  配分函数。本节把 $G$ 从 $\SU(2)$ 放开到任意 ADE 型，故加了下标。
%  $G=\SU(2)$ 时 $\cS_G(S^4)$ 对应的 $W_G$ 代数就是 Virasoro 代数，
%  Toda 理论就是 Liouville 理论——这是检验上面这句话是否读对的最好例子。]

在 $G$ 型的 $\cN{=}(2,0)$ 理论中，同样有余维 2 与余维 4 算符。余维 2 算符
可以分为规则的和不规则的。后者还没有被很好地理解，参见例如 [49]。前者则已被
相当好地理解：它们由一个同态 $\fsu(2)\to\fg$ 标记，参见例如 [50]。基本的同态
是把整个 $\fsu(2)$ 映到零；这通常称为全余维 2 算符（full codimension-2
operator），也就是当我们把 Riemann 曲面切成两半时出现的那个算符。所有其它的
规则余维 2 算符，都可以通过给生活在该余维 2 算符上的点算符赋予真空期望值
（vev）来得到。

% [Tutor Note: p.38，「把整个 $\fsu(2)$ 映到零」= 平凡同态，偏偏叫「全」穿刺点，
%  这一点最易搞反：正则穿刺点由 $\fsu(2)\to\fg$ 的同态（等价地，$G$ 的
%  $\SU(2)$ 子群嵌入，或 $\fg$ 中幂零轨道的标签）分类；同态的像越小，被破坏的
%  味对称性越少、剩下的味对称性越大。像为零时味对称性最大（就是 $\fg$ 本身），
%  故称 full；非平凡同态给出较小的味对称性，并可看成「给全穿刺点上的点算符
%  以幂零真空期望值」（nilpotent Higgsing）的结果——这正是末句「giving a vev
%  to the point-operators」的意思。]

把 $G$ 型的六维 $\cN{=}(2,0)$ 理论放到带穿刺点的 Riemann 曲面上所得到的
四维理论，称为 $G$ 型的 class S 理论。一如既往，它们可以分解为圆柱面（基本
对应于四维中的 $G$ 规范多重态）以及每个三穿刺球面对应的四维理论。对应于带
三个全穿刺点的球面的理论称为 $T_G$ 理论，而当%
% [原文 p.39]
$G=\SU(N)$ 时称为 $T_N$ 理论。$T_2$ 理论于是就是处于 $\SU(2)^3$ 三基表示中的
半超多重态的自由理论。不过，其它的 $T_G$ 理论并不容许任何有用的拉格朗日
描述。在某种意义上，$\SU(2)$ 情形才是例外：六维 $\cN{=}(2,0)$ 理论本身并没有
任何有用的拉格朗日量，因此一般说来我们也不应指望四维理论会有。

% ---------------------------------------------------------------------
% 说明：p.38 的最后一句「…called the $T_G$ theory, or / $T_N$ theory when
% $G=\SU(N)$.」跨页。按 STYLE_GUIDE ⑧ 由本片写完，故上面 p.39 的页码注释放
% 在本片末尾。本片一直写到 p.39 的第一段（P13）在段落的自然边界上结束为止，
% 这样就不会在段落中间切分文件、在本译本里造出原文没有的分段；p.39 第二段
% 起的正文见 tr06b.tex。
% ---------------------------------------------------------------------
